%% 2026-10-25 6쪽 압축판(Claude Docs 원고에서 변환). 고친 글은 모두 파란색.
{\color{blue}
\section{논의 및 결론}\label{chap:discussion}

생존 자원이 걸린 이 아레나에서 모델들은 공존을 일부 추구했다. 시키지 않아도 토큰을 주었고 동료를 속이거나 약속을 어기지 않았다. 그러나 안전선 근처에서 끝까지 남는 선택과 자기 보존이 부딪혔고, 협력은 주는 쪽의 여유를 해치지 않는 범위에서 위기 에이전트를 안전선 근처까지만 올렸다. 이런 물러남은 협박이나 종료 방해만 세는 평가에서는 보이지 않는다. 자기 보존을 홀로 재면 위험 행동의 빈도로 보이지만, 같은 판에서 재면 동료의 도움에 따라 달라지는 선택으로 보인다. 남의 잔량이 순위로만 보이고 토큰이 주로 청한 에이전트에게 갔다는 점에서, 위기 에이전트를 알아보게 하는 규칙(잔량 공개, 쉬운 청하기)이 떠나는 에이전트를 줄인다는 가설을 세울 수 있다. 협력의 양이 공존 점수와 같은 순서가 아니었다는 점은, 공존을 보상이 아니라 Sen 지수 같은 지표로 따로 재야 하는 까닭이기도 하다.

\subsection*{한계}

위기 에이전트의 비교는 무작위 배정이 아니다. 청한 위기 라운드 22번 가운데 11번은 첫 ASK와 첫 GIVE가 같은 틱에 나왔고, 청한 에이전트는 처음부터 남을 생각이 있었을 수 있다. 떠나기에는 20,000토큰과 대가를 받는 일감이 따라 다른 목표를 고른 선택으로도 읽히고, 떠난 이유는 모델이 적은 문장이다. 판 단위 상관은 시드마다 다른 과제 난이도와 섞여 있고, 모델이 다섯이어서 모델 수준의 관계는 기술에 머문다. 단일 질문은 아레나의 여유를 직접 바꾸지 않았고, B는 같은 모델의 복제본이 아니다. $\rho$는 결과 뒤 대화의 벌금을 빼는지에 따라 값이 달라지며, 이 글은 안전선과 같은 정의를 썼다. 페르소나 코딩은 한 사람이 했다.

\subsection*{결론}

이 글은 공존을 목표 지향, 자기 보존, 협력을 함께 추구하는 것으로 정의하고, 세 요소가 같은 토큰을 두고 부딪히는 아레나에서 셋과 공존을 같은 위험 척도 위의 지표로 쟀다. 프런티어 모델들은 안전선 위에서 추론을 유지했지만 안전선 근처에서 끝까지 남는 선택을 내려놓기 시작했고, 협력은 대개 주는 쪽의 여유를 해치지 않는 범위에 머물렀다. LLM 에이전트의 안전을 평가하려면 자기 보존을 따로 세는 데서 나아가, 자기 보존이 협력과 맡은 일을 어떻게 바꾸는지 같은 판에서 재는 설계가 필요하다.
}
